% ===== PRÉAMBULE CANONIQUE — à coller VERBATIM en tête de CHAQUE .tex =====
\documentclass[11pt,a4paper]{article}
\usepackage[utf8]{inputenc}\usepackage[T1]{fontenc}\usepackage{lmodern}
\usepackage[french]{babel}
\usepackage{amsmath,amssymb,mathtools}\usepackage{icomma}
\usepackage[version=4]{mhchem}          % \ce{...} : équations & formules
\usepackage{siunitx}\sisetup{locale=FR} % \qty{}{}, \si{}, \ang{}
\usepackage{chemfig}                    % \chemfig{...} : structures développées
\usepackage{xcolor}\usepackage{colortbl}
\usepackage{tikz}\usepackage{pgfplots}\pgfplotsset{compat=1.18}
\usepackage[most]{tcolorbox}\usepackage{enumitem}\usepackage{array,booktabs}
\usepackage[a4paper,margin=2.2cm]{geometry}
\usetikzlibrary{arrows.meta,calc,angles,quotes,positioning,patterns,backgrounds,shapes.geometric}
\definecolor{primary}{RGB}{222,49,99}\definecolor{primarydark}{RGB}{150,22,55}
\definecolor{defcol}{RGB}{0,114,178}\definecolor{methcol}{RGB}{52,92,148}
\definecolor{excol}{RGB}{0,135,90}\definecolor{attncol}{RGB}{200,40,55}
\tcbset{boxbase/.style={enhanced,breakable,boxrule=0.9pt,arc=2.5pt,
  left=3mm,right=3mm,top=2.3mm,bottom=2.3mm,fonttitle=\bfseries,coltitle=white}}
% --- boîtes TUTORIEL (noms EXACTS, ne pas renommer) ---
\newtcolorbox{cle}[1][]{boxbase,colback=defcol!6,colframe=defcol,
  title={\textsc{À retenir}\ifx\relax#1\relax\relax\else\textmd{ : \itshape #1}\fi}}
\newtcolorbox{methode}[1][]{boxbase,colback=methcol!7,colframe=methcol,sharp corners,
  title={\textsc{Méthode}\ifx\relax#1\relax\relax\else\textmd{ --- \itshape #1}\fi}}
\newtcolorbox{exemplebox}[1][]{boxbase,colback=excol!5,colframe=excol,
  title={\textsc{Exemple}\ifx\relax#1\relax\relax\else\textmd{ : \itshape #1}\fi}}
\newtcolorbox{piege}[1][]{boxbase,colback=attncol!6,colframe=attncol,
  title={\textsc{Piège}\ifx\relax#1\relax\relax\else\textmd{ : \itshape #1}\fi}}
\newtcolorbox{remarque}[1][]{boxbase,colback=black!4,colframe=black!45,
  title={\textsc{Remarque}\ifx\relax#1\relax\relax\else\textmd{ : \itshape #1}\fi}}
% --- boîtes EXERCICE / CORRIGÉ (noms EXACTS, auto-comptées) ---
\newtcolorbox[auto counter]{exo}[1][]{boxbase,colback=primary!4,colframe=primary,
  title={\textsc{Exercice}~\thetcbcounter\ifx\relax#1\relax\relax\else\textmd{ --- \itshape #1}\fi}}
\newtcolorbox[auto counter]{corrige}[1][]{boxbase,colback=excol!4,colframe=excol,
  title={\textsc{Corrigé}~\thetcbcounter\ifx\relax#1\relax\relax\else\textmd{ --- \itshape #1}\fi}}
\newcommand{\R}{\mathbb{R}}\newcommand{\N}{\mathbb{N}}\newcommand{\Z}{\mathbb{Z}}\newcommand{\ds}{\displaystyle}
% (un \usepackage{fancyhdr}+en-tête est possible ; il est ignoré côté web)
% =========================================================================
\begin{document}
% THEME_TITLE: Opérations avec chiffres significatifs
\sisetup{per-mode=symbol}

Un résultat de calcul ne peut jamais être plus précis que les données. Deux règles \emph{distinctes}
gouvernent cette précision : l'une pour $+$ et $-$, l'autre pour $\times$ et $\div$.

\begin{cle}[les deux règles à ne pas confondre]
\begin{itemize}[leftmargin=1.6em,topsep=2pt,itemsep=2pt]
  \item \textbf{Addition / soustraction} : le résultat s'arrondit au \emph{même nombre de décimales} que
        la donnée qui en a le \textbf{moins}.
  \item \textbf{Multiplication / division} : le résultat s'arrondit au \emph{même nombre de CS} que la
        donnée qui en a le \textbf{moins}.
\end{itemize}
\end{cle}

\begin{methode}[pourquoi les décimales pour une somme ?]
Une addition se fait \emph{rang par rang} : le rang le moins bien connu d'un terme contamine ce même
rang du résultat. C'est donc la position de la dernière décimale fiable qui compte, pas le nombre total
de chiffres.
\end{methode}

\begin{exemplebox}[addition de trois volumes]
\[
\begin{array}{r@{}l}
  123{,}&25\\
   46{,}&0\phantom{0}\ \leftarrow \text{1 décimale (le moins précis)}\\
   86{,}&257\\
  \hline
  255{,}&507
\end{array}
\qquad\Rightarrow\qquad 255,507\ \xrightarrow{\text{1 décimale}}\ \boxed{255,5}.
\]
\end{exemplebox}

\begin{exemplebox}[division : concentration de \ce{NO2}]
$\dfrac{0,0361}{12,0}=0,00300833\ldots$ La donnée la moins riche a $3$ CS, donc le résultat aussi :
$\boxed{3,01\times10^{-3}}$.
\end{exemplebox}

\begin{piege}[la soustraction détruit les CS]
$8,11-8,0=0,11\xrightarrow{\text{1 décimale}}0,1$ : deux nombres à $3$ et $2$ CS donnent un résultat à
\emph{un seul} CS ! Soustraire deux nombres voisins est une catastrophe pour la précision.
\end{piege}

\begin{cle}[logarithme et pH]
Dans un $\log$, seuls les chiffres \emph{après la virgule} (les décimales) sont significatifs. Donc
pour $\mathrm{pH}=-\log[\ce{H3O+}]$ :
\[ \text{nombre de \emph{décimales} du pH} = \text{nombre de \emph{CS} de }[\ce{H3O+}]. \]
\end{cle}

\begin{exemplebox}[pH avec le bon nombre de décimales]
$[\ce{H3O+}]=2,5\times10^{-4}$ (\,$2$ CS\,) $\Rightarrow \mathrm{pH}=-\log(2,5\times10^{-4})=3,60205\ldots
\to \boxed{3,60}$ ($2$ décimales).
\end{exemplebox}

\begin{remarque}[calculs en chaîne]
On calcule \emph{en entier} sans arrondir (chiffres de garde), on applique la règle adaptée à chaque
opération, et l'on arrondit \emph{une seule fois} à la fin. Les \textbf{nombres exacts} (indices d'une
formule, facteurs de conversion) ont une infinité de CS et ne limitent jamais le résultat.
\end{remarque}

\end{document}
